\documentclass[10pt,a4paper]{article}
\usepackage[margin=1.5cm]{geometry}
% common.tex - shared preamble for all ELL7282 A2 documents
\usepackage[T1]{fontenc}
\usepackage{lmodern}
\usepackage{amsmath,amssymb}
\usepackage{booktabs,longtable,tabularx,array,multirow}
\usepackage[table]{xcolor}
\usepackage{bytefield}
\usepackage{tikz}
\usetikzlibrary{arrows.meta,positioning,fit,calc,shapes.geometric}
\usepackage{pdflscape}
\usepackage{enumitem}
\usepackage{caption}
\usepackage[hidelinks]{hyperref}
\setlist{nosep,leftmargin=1.4em}
\setlength{\parskip}{3pt}
\setlength{\parindent}{0pt}

\newcommand{\project}{WM-8 five-stage pipelined washing-machine microcontroller}
\newcommand{\course}{ELL7282/ELL782 Introduction to Computer Architecture -- Assignment 2}
\newcommand{\pass}{\textcolor{green!45!black}{\textbf{PASS}}}
\newcommand{\fail}{\textcolor{red!75!black}{\textbf{FAIL}}}
% invalid (V=0) pipeline entry: grey with an empty-set marker
\newcommand{\bub}[1]{\textcolor{black!45}{$\varnothing$\,\texttt{#1}}}
\newcommand{\sig}[1]{\texttt{#1}}
\newcommand{\hex}[1]{\texttt{0x#1}}

% cycle-trace table (used by generated/trace*.tex)
\newenvironment{tracetab}{%
  \scriptsize\setlength{\tabcolsep}{2.5pt}\renewcommand{\arraystretch}{1.15}%
  \begin{longtable}{r p{3.0cm} p{3.0cm} p{3.0cm} p{3.0cm} p{3.0cm} c c c c p{4.6cm}}
  \toprule
  Cyc & IF (fetch PC) & ID (IF/ID) & EX (ID/EX) & MEM (EX/MEM) & WB (MEM/WB) &
  Stall & Flush & Fwd A & Fwd B & Architectural writes (commit at end of cycle)\\
  \midrule\endhead
  \bottomrule\endfoot}{%
  \end{longtable}}

% validation table (used by generated/val_g*.tex)
\newenvironment{valtab}{%
  \footnotesize\renewcommand{\arraystretch}{1.2}%
  \begin{longtable}{|p{0.6cm}|p{3.7cm}|p{4.2cm}|p{5.6cm}|p{1.0cm}|}
  \hline
  \rowcolor{black!8}\textbf{ID} & \textbf{Initial inputs / state, stimulus} & \textbf{Expected result} &
  \textbf{Observed result} & \textbf{P/F}\\ \hline\endhead}{%
  \end{longtable}}

% program listing (used by generated/listing_*.tex)
\newenvironment{listtab}{%
  \footnotesize\setlength{\tabcolsep}{4pt}\renewcommand{\arraystretch}{1.05}%
  \begin{longtable}{l l c l l p{6.2cm}}
  \toprule
  Addr & Hex & Fmt & Label & Assembly & Comment\\
  \midrule\endhead
  \bottomrule\endfoot}{%
  \end{longtable}}

\begin{document}
\begin{center}
{\Large\textbf{Validation Report (Section 11)}}\\[2pt]
\project\ -- \course
\end{center}

\section{Method and reproducibility}
\begin{itemize}
\item \textbf{Device under test.} The six test benches \texttt{tests/t1}--\texttt{t6\_*.circ} each contain \emph{the same}
  sub-circuits as \texttt{washing\_pipeline.circ} (\texttt{cpu}, \texttt{mmio}, \texttt{timer8}, \texttt{safety} and all gate-level
  blocks). Only the top level differs. Its pieces are:
  \begin{itemize}
  \item the ROM holds the test program;
  \item a power-on-reset flip-flop asserts \texttt{RESET} for cycle 0;
  \item a 16-bit cycle counter drives a \emph{stimulus ROM} that supplies \texttt{START}, \texttt{DOOR\_CLOSED} and
    \texttt{MODE} (the sensor values are listed below);
  \item a small plant model supplies \texttt{WATER\_LEVEL} and \texttt{DRUM\_RPM}: water $+3$ every 16 cycles while
    \texttt{VALVE\_IN}=1 and $-7$ (saturating at 0) every 16 cycles while \texttt{PUMP\_OUT}=1; drum 200~rpm while
    \texttt{MOTOR\_EN\_ACTUAL}=1, otherwise $-40$ per cycle down to 0;
  \item \texttt{TEMP} is the constant \hex{19};
  \item an output pin \texttt{halt} stops the run after a fixed number of cycles.
  \end{itemize}
  Inputs change only right after a rising edge, so they are stable at every sampling edge.
\item \textbf{Observation.} One row per clock cycle with 36 monitored signals: pipeline state, stall/flush, forwarding,
  write-back and store ports, \texttt{RAM\_WE}, actuator latch, display, gated motor output, \texttt{TRIP}, timer
  registers and sensor values. The \emph{observed} column below was extracted automatically by
  \texttt{tools/make\_reports.py} from these rows. The rows were recorded by cycle-accurate simulation of the
  test-bench \texttt{.circ} netlists with \texttt{tools/circsim.py}. That simulator rebuilds connectivity from the XML
  geometry and also checks for width conflicts, multiple drivers, undriven nets and combinational loops (none found).
\item \textbf{Reproducing in Logisim.} \texttt{java -jar logisim-generic-2.7.1.jar tests/t1\_hazards.circ -tty table > t1.txt},
  then \texttt{diff t1.txt tests/expected/t1\_hazards.txt}; \texttt{tests/run\_logisim\_tests.sh} does this for all six. The same
  files can also be opened in the GUI and clocked by hand (\texttt{RUNNING\_IN\_LOGISIM.md}).
\item \textbf{Cross-checks.} For t1--t3 the ordered list of register write-backs and stores is compared with an
  instruction-level reference model (\texttt{tools/isa\_ref.py}). In a strict mode in which every register and RAM byte starts
  \emph{undefined}, the reference model also confirms that \texttt{program.asm} never reads uninitialised state
  (normal path and safety-trip path).
\end{itemize}

\textbf{Result: 39 of 39 checks passed.}

\section*{Sensor inputs used}
\begin{center}\footnotesize
\begin{tabular}{l l}
\toprule
Bench & Inputs (cycle numbers are rising-edge counts; values hold until the next change)\\
\midrule
t1--t4 & \texttt{DOOR\_CLOSED}=1, \texttt{START}=0, \texttt{MODE}=00, \texttt{WATER\_LEVEL}=0, \texttt{TEMP}=\hex{19}, \texttt{DRUM\_RPM}=0 (not used by the programs)\\
t5 & \texttt{DOOR}: 1, 0 @40, 1 @52, 0 @88, 1 @108; \texttt{START}: 0, 1 @64, 0 @76, 1 @92, 0 @100, 1 @120, 0 @140\\
t6 & \texttt{DOOR}=1; \texttt{START}: 1 @32, 0 @128, 1 @1600, 0 @1696; \texttt{MODE}: 00, 01 @480 (during run 1), 01 @1600, 00 @2400 (during run 2);\\
   & \texttt{WATER\_LEVEL}/\texttt{DRUM\_RPM} from the plant model (fill reads \dots 96, 99, 102 and \dots 156, 159, 162)\\
\bottomrule
\end{tabular}
\end{center}

\section{Group 1 -- ALU forwarding, load-use stall, youngest-writer priority (\texttt{t1\_hazards})}
Initial state: all pipeline entries invalid after reset. Registers are initialised by the program
(R1=01, R2=10, R3=03, R4=20, RAM[10]=03). Cycle traces are in \texttt{pipeline\_traces.pdf}.
\begin{valtab}
1a & \S8.1 \texttt{ADD R1,R2; SUB R3,R1}; R1=01, R2=10, R3=03 & no stall; SUB takes R1 from EX/MEM; R3=F2 & cyc 10: stall=0, FwdB=EX/MEM; R3=F2 (cyc 12) & \pass \\ \hline
1b & \S8.2 \texttt{LD R1,[R2]; ADD R3,R1}; RAM[10]=03 & exactly 1 stall cycle; PC and IF/ID hold; bubble (V=0) in EX; R1 from MEM/WB; R3=F5 & stall cycles [11]; cyc 12: IDEX.V=0, PC held at 0A; cyc 13: FwdB=MEM/WB; R3=F5 & \pass \\ \hline
1c & \S8.3 \texttt{ADD R1,R2; SUB R1,R3; ST R1,[R4]} & store data from youngest writer (SUB, EX/MEM) = 1E, not older ADD (MEM/WB) = 13 & cyc 16: FwdA=EX/MEM while MEM/WB holds pc 0A; RAM[20]$\leftarrow$1E (cyc 17) & \pass \\ \hline
1d & read-back \texttt{LD R5,[R4]; ST R5,[DISPLAY]; HALT; LDI R6,FF} & DISPLAY=1E; HALT stops machine; younger LDI R6 never writes & DISPLAY=1E; HALTED=1; writes to R6: 0 & \pass \\ \hline
1e & whole program vs. ISA reference model & identical ordered list of register/memory writes & 15 writes, identical & \pass \\ \hline

\end{valtab}

\section{Group 2 -- taken branch kills a wrong-path actuator store and a wrong-path HALT (\texttt{t2\_branch\_halt})}
Initial state: reset. R2=\hex{21} (VALVE+LOCK pattern), R4=\hex{F4} (actuator latch), R3=07, R1=09.
\begin{valtab}
2a & \S8.4 \texttt{SUB R1,R1; BEQZ R1,T; ST R2,[Rvalve]; ADD R3,R4}; R2=21, Rvalve=F4 & branch taken in EX (R1 via EX/MEM); IF/ID and ID/EX invalidated; wrong-path store reaches MEM with V=0; VALVE/ACT never written & cyc 9: flush=1, FwdA=EX/MEM; cyc 10: IFID.V=0, IDEX.V=0; cyc 11: ST in MEM V=0; valid stores to F4: 0; ACT=00 throughout & \pass \\ \hline
2b & wrong-path \texttt{ADD R3,R4} & R3 keeps 07 (DISPLAY shows 07) & writes to R3: 07 (only the LDI); DISPLAY=07 at cyc 18 & \pass \\ \hline
2c & \texttt{BNEZ R1,bad} with R1=0 & not taken, no flush & cyc 12: flush=0 & \pass \\ \hline
2d & wrong-path \texttt{HALT} at 0B (after taken BNEZ) and 0F/10 (after JMP) & flushed before EX; execution continues to DISPLAY=AA & wrong-path HALTs valid in EX: 0; DISPLAY=AA at cyc 23; HALTED first at cyc 23 by HALT @13 & \pass \\ \hline
2e & real \texttt{HALT} followed by \texttt{ST R2,[R4]} & younger store killed; fetch stops; remains halted & store @14 valid in MEM: False; ACT=00; PC frozen at 15 & \pass \\ \hline
2f & whole program vs. ISA reference model & identical write sequence & 9 writes, identical & \pass \\ \hline

\end{valtab}

\section{Group 3 -- RAM store/load and reserved-address isolation (\texttt{t3\_ram\_reserved})}
Initial state: reset. The program first writes RAM[00]=5A, RAM[6F]=33 and RAM[7F]=A5, then stores \hex{77} to addresses that would alias
these locations if bit 7 were ignored (\hex{80}$\rightarrow$00, \hex{EF}$\rightarrow$6F, \hex{FF}$\rightarrow$7F) and to read-only
\hex{F0}.
\begin{valtab}
3a & \texttt{ST} to RAM 00, 6F, 7F & RAM write-enable exactly 3 times & RAM\_WE at cyc 6 [00]$\leftarrow$5A, cyc 9 [6F]$\leftarrow$33, cyc 12 [7F]$\leftarrow$A5 & \pass \\ \hline
3b & \texttt{ST 77} to 80, EF (reserved), FF (unimpl.), F0 (read-only) & no RAM write (no aliasing), no latch change & 4 stores seen on the bus, RAM\_WE during them: 0; ACT/DISPLAY unchanged & \pass \\ \hline
3c & loads from 7F, 00, 6F after the reserved stores & R6=A5, R7=5A, R1=33 (unchanged) & R6=A5, R7=5A, R1=33 & \pass \\ \hline
3d & loads from 80, EF (reserved) and F7 (unimplemented) & read as 00 & R0=00, R3=00, R5=00 & \pass \\ \hline
3e & whole program vs. ISA reference model & identical write sequence & 31 writes, identical & \pass \\ \hline

\end{valtab}

\section{Group 4 -- timer (\texttt{t4\_timer})}
Initial state: reset (\texttt{COUNT}=\texttt{LIMIT}=0, \texttt{EN}=\texttt{MATCH}=0).
\begin{valtab}
4a & LIMIT=6, CTRL=03 (clear+enable) & COUNT=0 after the CTRL edge, +1 per clock; MATCH set on the edge where COUNT becomes 6 & CTRL write cyc 11; COUNT=00 at cyc 12; MATCH=1 at cyc 18 with COUNT=06 & \pass \\ \hline
4b & polling loops \texttt{LD R2,[T\_STATUS]; BEQZ R2} (load-use stall) & COUNT increments during every stall cycle & stall cycles [16, 21, 31, 36, 41, 53]: COUNT +1 each & \pass \\ \hline
4c & keep running after MATCH & MATCH stays 1 (sticky) while COUNT passes LIMIT & MATCH=1 from cyc 18 to 30, COUNT reached 12 & \pass \\ \hline
4d & re-arm: LIMIT=0A, CTRL=03 & clear wins: COUNT=0, MATCH=0 on the CTRL edge; MATCH again at COUNT=0A & cyc 31: COUNT=00, MATCH=0; MATCH at cyc 41 (COUNT=0A) & \pass \\ \hline
4e & timer stopped (CTRL=00), then LIMIT := current COUNT & no retroactive MATCH & cyc 57: COUNT=05, LIMIT=05, MATCH=0; DISPLAY(MATCH)=00 & \pass \\ \hline
4f & LIMIT=03, CTRL=01 (enable, no clear), then \texttt{HALT} & timer keeps ticking while halted; COUNT wraps FF$\rightarrow$00; MATCH at COUNT=03 & wrap at cyc 315 (HALTED=1); MATCH at cyc 318 (COUNT=03) & \pass \\ \hline

\end{valtab}
\begin{center}\footnotesize
\begin{tabular}{rcccccl}\toprule Cyc & STALL & HALTED & EN & COUNT & LIMIT & MATCH \\ \midrule
10 & 0 & 0 & 0 & 00 & 06 & 0 \\
11 & 1 & 0 & 0 & 00 & 06 & 0 \\
12 & 0 & 0 & 1 & 00 & 06 & 0 \\
13 & 0 & 0 & 1 & 01 & 06 & 0 \\
14 & 0 & 0 & 1 & 02 & 06 & 0 \\
15 & 0 & 0 & 1 & 03 & 06 & 0 \\
16 & 1 & 0 & 1 & 04 & 06 & 0 \\
17 & 0 & 0 & 1 & 05 & 06 & 0 \\
18 & 0 & 0 & 1 & 06 & 06 & 1 \\
19 & 0 & 0 & 1 & 07 & 06 & 1 \\
\midrule \multicolumn{7}{c}{\dots} \\ \midrule
313 & 0 & 1 & 1 & FE & 03 & 0 \\
314 & 0 & 1 & 1 & FF & 03 & 0 \\
315 & 0 & 1 & 1 & 00 & 03 & 0 \\
316 & 0 & 1 & 1 & 01 & 03 & 0 \\
317 & 0 & 1 & 1 & 02 & 03 & 0 \\
318 & 0 & 1 & 1 & 03 & 03 & 1 \\
319 & 0 & 1 & 1 & 04 & 03 & 1 \\
\bottomrule\end{tabular}
\\[2pt]
Excerpt: interval 1 (stall cycles marked by \texttt{STALL}=1, the count still advances) and the wrap-around while halted.
\end{center}

\section{Group 5 -- safety interlock (\texttt{t5\_safety})}
Initial state: reset (\texttt{TRIP}=0, \texttt{START\_PREV}=0). The program requests \texttt{MOTOR\_EN} and keeps the request
until it sees a START press and release, then clears the request and waits for \texttt{TRIP}=0 before requesting
the motor again. The physical output \texttt{MOTOR\_EN\_ACTUAL} is checked separately from the CPU request
(\texttt{ACT} bit 2).
\begin{valtab}
5a & DOOR=1, software writes ACT=24 (LOCK+MOTOR\_EN) & ActualMotorEnable=1 & cyc 11: ACT=24, actual=1 & \pass \\ \hline
5b & open door at cyc 40 (request still 1) & actual motor drops in the same cycle (combinational); TRIP latched at next edge & cyc 40: MOTOR\_EN req=1, actual=0, TRIP=0; cyc 41: TRIP=1 & \pass \\ \hline
5c & re-close door at cyc 52, request still 1 & no automatic restart & cyc 52-61: actual=0, TRIP=1, request=1 & \pass \\ \hline
5d & START 0$\rightarrow$1 at cyc 64 while MOTOR\_EN request still active & not a valid re-arm: TRIP stays 1 & cyc 65: TRIP=1 & \pass \\ \hline
5e & software clears request (cyc 83); START edge with door open (cyc 92) & not a valid re-arm: TRIP stays 1 & cyc 93: TRIP=1 & \pass \\ \hline
5f & door closed, request off, new START edge at cyc 120 & TRIP cleared; motor runs only after software requests it again & cyc 121: TRIP=0; actual=1 again at cyc 127 (ACT=24); DISPLAY=03 & \pass \\ \hline

\end{valtab}
\begin{center}\footnotesize
\begin{tabular}{rccccccl}\toprule Cyc & DOOR & START & ACT & MOTOR\_EN req & Actual & TRIP & note \\ \midrule
0 & 1 & 0 & 00 & 0 & 0 & 0 & \\
11 & 1 & 0 & 24 & 1 & 1 & 0 & \\
40 & 0 & 0 & 24 & 1 & 0 & 0 & \\
41 & 0 & 0 & 24 & 1 & 0 & 1 & \\
52 & 1 & 0 & 24 & 1 & 0 & 1 & \\
64 & 1 & 1 & 24 & 1 & 0 & 1 & \\
76 & 1 & 0 & 24 & 1 & 0 & 1 & \\
83 & 1 & 0 & 20 & 0 & 0 & 1 & \\
88 & 0 & 0 & 20 & 0 & 0 & 1 & \\
92 & 0 & 1 & 20 & 0 & 0 & 1 & \\
100 & 0 & 0 & 20 & 0 & 0 & 1 & \\
108 & 1 & 0 & 20 & 0 & 0 & 1 & \\
120 & 1 & 1 & 20 & 0 & 0 & 1 & \\
121 & 1 & 1 & 20 & 0 & 0 & 0 & \\
127 & 1 & 1 & 24 & 1 & 1 & 0 & \\
140 & 1 & 0 & 24 & 1 & 1 & 0 & \\
\bottomrule\end{tabular}
\\[2pt]
Every change of \texttt{DOOR}, \texttt{START}, \texttt{ACT}, actual motor enable or \texttt{TRIP}.
\end{center}

\section{Group 6 -- end-to-end wash, two MODE settings (\texttt{t6\_wash\_e2e}, \texttt{program.asm})}
Initial state: reset, \texttt{DOOR\_CLOSED}=1, empty drum. Run 1 uses \texttt{MODE}=00 (NORMAL: target 100, wash 64, spin 32,
drain $\le$10). Run 2 starts after run 1 has completed and the buzzer is on, with \texttt{MODE}=01 (HEAVY: 160, 96, 48, $\le$10).
\texttt{MODE} is changed during both runs to show that the snapshot is used. The overshoot samples are 99$\rightarrow$102
(target 100) and 159$\rightarrow$162 (target 160). The drain thresholds are also crossed by a jump (11$\rightarrow$4 and 15$\rightarrow$8).
\begin{valtab}
6a & NORMAL (MODE=00): START=1, DOOR=1; plant +3 per 16 cycles & fill until WATER\_LEVEL $\geq$ 100 (unsigned) although reading jumps past 100 & last samples 99, 99, 99, 102; valve closed at cyc 619 with WATER\_LEVEL=102 & \pass \\ \hline
6b & wash phase & LIMIT=64; motor direction alternates; motor off at MATCH & LIMIT=64; DIR changed 3 times; motor off at cyc 710 (COUNT=75, MATCH=1) & \pass \\ \hline
6c & drain phase (plant $-7$ per 16 cycles) & pump until WATER\_LEVEL $\leq$ 10 (unsigned) & samples \dots 11, 11, 4; spin starts cyc 950 & \pass \\ \hline
6d & spin phase & timer cleared/re-armed, LIMIT=32, MATCH at 32 & LIMIT=32, COUNT=2 at spin start, MATCH at cyc 980 (COUNT=32) & \pass \\ \hline
6e & completion & all actuators off; unlock only after drum RPM=0; BUZZER=1; DISPLAY=08 & motor/pump off cyc 989; RPM=0 cyc 994; ACT=40 (unlock+buzzer) cyc 1004; DISPLAY=08 & \pass \\ \hline
6g & HEAVY (MODE=01): START=1, DOOR=1; plant +3 per 16 cycles & fill until WATER\_LEVEL $\geq$ 160 (unsigned) although reading jumps past 160 & last samples 159, 159, 159, 162; valve closed at cyc 2507 with WATER\_LEVEL=162 & \pass \\ \hline
6h & wash phase & LIMIT=96; motor direction alternates; motor off at MATCH & LIMIT=96; DIR changed 5 times; motor off at cyc 2626 (COUNT=103, MATCH=1) & \pass \\ \hline
6i & drain phase (plant $-7$ per 16 cycles) & pump until WATER\_LEVEL $\leq$ 10 (unsigned) & samples \dots 15, 15, 8; spin starts cyc 2998 & \pass \\ \hline
6j & spin phase & timer cleared/re-armed, LIMIT=48, MATCH at 48 & LIMIT=48, COUNT=2 at spin start, MATCH at cyc 3044 (COUNT=48) & \pass \\ \hline
6k & completion & all actuators off; unlock only after drum RPM=0; BUZZER=1; DISPLAY=18 & motor/pump off cyc 3055; RPM=0 cyc 3060; ACT=40 (unlock+buzzer) cyc 3070; DISPLAY=18 & \pass \\ \hline
6m & MODE changed during run 1 (cyc 480) and run 2 (cyc 2400) & parameters of the running cycle unchanged & run 1 still used target 100 / LIMIT 64 / 32; run 2 still used 160 / 96 / 48 (rows 6a--6k) & \pass \\ \hline

\end{valtab}
\begin{center}\footnotesize
\begin{tabular}{lrr}\toprule Event & NORMAL & HEAVY \\ \midrule
DOOR\_LOCK commanded (START accepted) & 67 & 1631 \\
VALVE\_IN on & 78 & 1642 \\
VALVE\_IN off & 619 & 2507 \\
wash motor off (MATCH) & 710 & 2626 \\
spin start & 950 & 2998 \\
spin MATCH & 980 & 3044 \\
motor/pump off & 989 & 3055 \\
drum stopped (RPM=0) & 994 & 3060 \\
unlock + BUZZER & 1004 & 3070 \\
water level at valve close & 102 & 162 \\
last drain sample & 4 & 8 \\
\bottomrule\end{tabular}
\\[2pt]
Phase timeline (cycle numbers) of both runs. DISPLAY sequence NORMAL 01, 02, 03, 04, 05, 06, 07, 08; HEAVY 12, 13, 14, 15, 16, 17, 18.
\end{center}

\appendix
\section{Remaining test program listings}
\textbf{t3\_ram\_reserved.asm}
\begin{listtab}
\texttt{00} & \texttt{125A} & I & \texttt{} & \texttt{LDI  R1, 0x5A} & {\footnotesize } \\
\texttt{01} & \texttt{1400} & I & \texttt{} & \texttt{LDI  R2, 0x00} & {\footnotesize } \\
\texttt{02} & \texttt{7280} & R & \texttt{} & \texttt{ST   R1, [R2]} & {\footnotesize RAM[0x00] = 0x5A} \\
\texttt{03} & \texttt{1633} & I & \texttt{} & \texttt{LDI  R3, 0x33} & {\footnotesize } \\
\texttt{04} & \texttt{146F} & I & \texttt{} & \texttt{LDI  R2, 0x6F} & {\footnotesize } \\
\texttt{05} & \texttt{7680} & R & \texttt{} & \texttt{ST   R3, [R2]} & {\footnotesize RAM[0x6F] = 0x33} \\
\texttt{06} & \texttt{16A5} & I & \texttt{} & \texttt{LDI  R3, 0xA5} & {\footnotesize } \\
\texttt{07} & \texttt{147F} & I & \texttt{} & \texttt{LDI  R2, 0x7F} & {\footnotesize } \\
\texttt{08} & \texttt{7680} & R & \texttt{} & \texttt{ST   R3, [R2]} & {\footnotesize RAM[0x7F] = 0xA5} \\
\texttt{09} & \texttt{1A77} & I & \texttt{} & \texttt{LDI  R5, 0x77} & {\footnotesize } \\
\texttt{0A} & \texttt{1880} & I & \texttt{} & \texttt{LDI  R4, 0x80} & {\footnotesize } \\
\texttt{0B} & \texttt{7B00} & R & \texttt{} & \texttt{ST   R5, [R4]} & {\footnotesize reserved 0x80: ignored (would alias RAM[0x00])} \\
\texttt{0C} & \texttt{18EF} & I & \texttt{} & \texttt{LDI  R4, 0xEF} & {\footnotesize } \\
\texttt{0D} & \texttt{7B00} & R & \texttt{} & \texttt{ST   R5, [R4]} & {\footnotesize reserved 0xEF: ignored (would alias RAM[0x6F])} \\
\texttt{0E} & \texttt{18FF} & I & \texttt{} & \texttt{LDI  R4, 0xFF} & {\footnotesize } \\
\texttt{0F} & \texttt{7B00} & R & \texttt{} & \texttt{ST   R5, [R4]} & {\footnotesize unimplemented peripheral: ignored (alias RAM[0x7F])} \\
\texttt{10} & \texttt{18F0} & I & \texttt{} & \texttt{LDI  R4, 0xF0} & {\footnotesize } \\
\texttt{11} & \texttt{7B00} & R & \texttt{} & \texttt{ST   R5, [R4]} & {\footnotesize read-only IN\_STATUS: ignored} \\
\texttt{12} & \texttt{6C80} & R & \texttt{} & \texttt{LD   R6, [R2]} & {\footnotesize R6 = RAM[0x7F] = 0xA5} \\
\texttt{13} & \texttt{1400} & I & \texttt{} & \texttt{LDI  R2, 0x00} & {\footnotesize } \\
\texttt{14} & \texttt{6E80} & R & \texttt{} & \texttt{LD   R7, [R2]} & {\footnotesize R7 = RAM[0x00] = 0x5A} \\
\texttt{15} & \texttt{146F} & I & \texttt{} & \texttt{LDI  R2, 0x6F} & {\footnotesize } \\
\texttt{16} & \texttt{6280} & R & \texttt{} & \texttt{LD   R1, [R2]} & {\footnotesize R1 = RAM[0x6F] = 0x33} \\
\texttt{17} & \texttt{1880} & I & \texttt{} & \texttt{LDI  R4, 0x80} & {\footnotesize } \\
\texttt{18} & \texttt{6100} & R & \texttt{} & \texttt{LD   R0, [R4]} & {\footnotesize R0 = 0x00 (reserved reads 0)} \\
\texttt{19} & \texttt{18EF} & I & \texttt{} & \texttt{LDI  R4, 0xEF} & {\footnotesize } \\
\texttt{1A} & \texttt{6700} & R & \texttt{} & \texttt{LD   R3, [R4]} & {\footnotesize R3 = 0x00} \\
\texttt{1B} & \texttt{18F7} & I & \texttt{} & \texttt{LDI  R4, 0xF7} & {\footnotesize } \\
\texttt{1C} & \texttt{6B00} & R & \texttt{} & \texttt{LD   R5, [R4]} & {\footnotesize R5 = 0x00 (unimplemented peripheral)} \\
\texttt{1D} & \texttt{18F5} & I & \texttt{} & \texttt{LDI  R4, 0xF5} & {\footnotesize } \\
\texttt{1E} & \texttt{7D00} & R & \texttt{} & \texttt{ST   R6, [R4]} & {\footnotesize DISPLAY = 0xA5} \\
\texttt{1F} & \texttt{F000} & N & \texttt{} & \texttt{HALT} & {\footnotesize } \\

\end{listtab}
\textbf{t4\_timer.asm}
\begin{listtab}
\texttt{00} & \texttt{1EF9} & I & \texttt{} & \texttt{LDI  R7, 0xF9} & {\footnotesize T\_LIMIT} \\
\texttt{01} & \texttt{1CFA} & I & \texttt{} & \texttt{LDI  R6, 0xFA} & {\footnotesize T\_CTRL} \\
\texttt{02} & \texttt{1AFB} & I & \texttt{} & \texttt{LDI  R5, 0xFB} & {\footnotesize T\_STATUS} \\
\texttt{03} & \texttt{18F8} & I & \texttt{} & \texttt{LDI  R4, 0xF8} & {\footnotesize T\_COUNT} \\
\texttt{04} & \texttt{1206} & I & \texttt{} & \texttt{LDI  R1, 6} & {\footnotesize } \\
\texttt{05} & \texttt{73C0} & R & \texttt{} & \texttt{ST   R1, [R7]} & {\footnotesize LIMIT = 6} \\
\texttt{06} & \texttt{1203} & I & \texttt{} & \texttt{LDI  R1, 0x03} & {\footnotesize } \\
\texttt{07} & \texttt{7380} & R & \texttt{} & \texttt{ST   R1, [R6]} & {\footnotesize interval 1: clear + enable} \\
\texttt{08} & \texttt{6540} & R & \texttt{poll1} & \texttt{LD   R2, [R5]} & {\footnotesize } \\
\texttt{09} & \texttt{9408} & I & \texttt{} & \texttt{BEQZ R2, poll1} & {\footnotesize load-use stall every iteration (timer keeps counting)} \\
\texttt{0A} & \texttt{6700} & R & \texttt{} & \texttt{LD   R3, [R4]} & {\footnotesize COUNT kept counting past LIMIT} \\
\texttt{0B} & \texttt{6540} & R & \texttt{} & \texttt{LD   R2, [R5]} & {\footnotesize MATCH still 1 (sticky)} \\
\texttt{0C} & \texttt{120A} & I & \texttt{} & \texttt{LDI  R1, 10} & {\footnotesize } \\
\texttt{0D} & \texttt{73C0} & R & \texttt{} & \texttt{ST   R1, [R7]} & {\footnotesize LIMIT = 10} \\
\texttt{0E} & \texttt{1203} & I & \texttt{} & \texttt{LDI  R1, 0x03} & {\footnotesize } \\
\texttt{0F} & \texttt{7380} & R & \texttt{} & \texttt{ST   R1, [R6]} & {\footnotesize interval 2: clear COUNT+MATCH, re-arm} \\
\texttt{10} & \texttt{6540} & R & \texttt{} & \texttt{LD   R2, [R5]} & {\footnotesize MATCH = 0 after clear} \\
\texttt{11} & \texttt{6540} & R & \texttt{poll2} & \texttt{LD   R2, [R5]} & {\footnotesize } \\
\texttt{12} & \texttt{9411} & I & \texttt{} & \texttt{BEQZ R2, poll2} & {\footnotesize } \\
\texttt{13} & \texttt{1200} & I & \texttt{} & \texttt{LDI  R1, 0x00} & {\footnotesize } \\
\texttt{14} & \texttt{73C0} & R & \texttt{} & \texttt{ST   R1, [R7]} & {\footnotesize LIMIT = 0 (comparison disabled)} \\
\texttt{15} & \texttt{1203} & I & \texttt{} & \texttt{LDI  R1, 0x03} & {\footnotesize } \\
\texttt{16} & \texttt{7380} & R & \texttt{} & \texttt{ST   R1, [R6]} & {\footnotesize clear + enable} \\
\texttt{17} & \texttt{0000} & N & \texttt{} & \texttt{NOP} & {\footnotesize } \\
\texttt{18} & \texttt{0000} & N & \texttt{} & \texttt{NOP} & {\footnotesize } \\
\texttt{19} & \texttt{0000} & N & \texttt{} & \texttt{NOP} & {\footnotesize } \\
\texttt{1A} & \texttt{1200} & I & \texttt{} & \texttt{LDI  R1, 0x00} & {\footnotesize } \\
\texttt{1B} & \texttt{7380} & R & \texttt{} & \texttt{ST   R1, [R6]} & {\footnotesize disable (no clear): COUNT frozen} \\
\texttt{1C} & \texttt{6700} & R & \texttt{} & \texttt{LD   R3, [R4]} & {\footnotesize R3 = frozen COUNT} \\
\texttt{1D} & \texttt{77C0} & R & \texttt{} & \texttt{ST   R3, [R7]} & {\footnotesize LIMIT = COUNT -$>$ MATCH must stay 0} \\
\texttt{1E} & \texttt{6540} & R & \texttt{} & \texttt{LD   R2, [R5]} & {\footnotesize R2 = 0} \\
\texttt{1F} & \texttt{12F5} & I & \texttt{} & \texttt{LDI  R1, 0xF5} & {\footnotesize } \\
\texttt{20} & \texttt{7440} & R & \texttt{} & \texttt{ST   R2, [R1]} & {\footnotesize DISPLAY = MATCH = 0} \\
\texttt{21} & \texttt{1203} & I & \texttt{} & \texttt{LDI  R1, 0x03} & {\footnotesize } \\
\texttt{22} & \texttt{73C0} & R & \texttt{} & \texttt{ST   R1, [R7]} & {\footnotesize LIMIT = 3} \\
\texttt{23} & \texttt{1201} & I & \texttt{} & \texttt{LDI  R1, 0x01} & {\footnotesize } \\
\texttt{24} & \texttt{7380} & R & \texttt{} & \texttt{ST   R1, [R6]} & {\footnotesize enable, no clear: COUNT runs ..FF -$>$ 00 -$>$ 03 =$>$ MATCH} \\
\texttt{25} & \texttt{F000} & N & \texttt{} & \texttt{HALT} & {\footnotesize timer keeps ticking while halted} \\

\end{listtab}
\textbf{t5\_safety.asm}
\begin{listtab}
\texttt{00} & \texttt{1EF4} & I & \texttt{} & \texttt{LDI  R7, 0xF4} & {\footnotesize ACT} \\
\texttt{01} & \texttt{1CF0} & I & \texttt{} & \texttt{LDI  R6, 0xF0} & {\footnotesize IN\_STATUS} \\
\texttt{02} & \texttt{1AF5} & I & \texttt{} & \texttt{LDI  R5, 0xF5} & {\footnotesize DISPLAY} \\
\texttt{03} & \texttt{1801} & I & \texttt{} & \texttt{LDI  R4, 0x01} & {\footnotesize START mask} \\
\texttt{04} & \texttt{1610} & I & \texttt{} & \texttt{LDI  R3, 0x10} & {\footnotesize TRIP mask} \\
\texttt{05} & \texttt{1224} & I & \texttt{} & \texttt{LDI  R1, 0x24} & {\footnotesize } \\
\texttt{06} & \texttt{73C0} & R & \texttt{} & \texttt{ST   R1, [R7]} & {\footnotesize request DOOR\_LOCK + MOTOR\_EN} \\
\texttt{07} & \texttt{1201} & I & \texttt{} & \texttt{LDI  R1, 0x01} & {\footnotesize } \\
\texttt{08} & \texttt{7340} & R & \texttt{} & \texttt{ST   R1, [R5]} & {\footnotesize DISPLAY = 1 : motor requested} \\
\texttt{09} & \texttt{6580} & R & \texttt{waitS} & \texttt{LD   R2, [R6]} & {\footnotesize } \\
\texttt{0A} & \texttt{5500} & R & \texttt{} & \texttt{AND  R2, R4} & {\footnotesize } \\
\texttt{0B} & \texttt{9409} & I & \texttt{} & \texttt{BEQZ R2, waitS} & {\footnotesize wait for a START press (request still active)} \\
\texttt{0C} & \texttt{6580} & R & \texttt{waitR} & \texttt{LD   R2, [R6]} & {\footnotesize } \\
\texttt{0D} & \texttt{5500} & R & \texttt{} & \texttt{AND  R2, R4} & {\footnotesize } \\
\texttt{0E} & \texttt{A40C} & I & \texttt{} & \texttt{BNEZ R2, waitR} & {\footnotesize wait for START release} \\
\texttt{0F} & \texttt{1220} & I & \texttt{} & \texttt{LDI  R1, 0x20} & {\footnotesize } \\
\texttt{10} & \texttt{73C0} & R & \texttt{} & \texttt{ST   R1, [R7]} & {\footnotesize software clears MOTOR\_EN request} \\
\texttt{11} & \texttt{1202} & I & \texttt{} & \texttt{LDI  R1, 0x02} & {\footnotesize } \\
\texttt{12} & \texttt{7340} & R & \texttt{} & \texttt{ST   R1, [R5]} & {\footnotesize DISPLAY = 2 : waiting for re-arm} \\
\texttt{13} & \texttt{6580} & R & \texttt{waitT} & \texttt{LD   R2, [R6]} & {\footnotesize } \\
\texttt{14} & \texttt{54C0} & R & \texttt{} & \texttt{AND  R2, R3} & {\footnotesize } \\
\texttt{15} & \texttt{A413} & I & \texttt{} & \texttt{BNEZ R2, waitT} & {\footnotesize until hardware TRIP is cleared} \\
\texttt{16} & \texttt{1224} & I & \texttt{} & \texttt{LDI  R1, 0x24} & {\footnotesize } \\
\texttt{17} & \texttt{73C0} & R & \texttt{} & \texttt{ST   R1, [R7]} & {\footnotesize request motor again} \\
\texttt{18} & \texttt{1203} & I & \texttt{} & \texttt{LDI  R1, 0x03} & {\footnotesize } \\
\texttt{19} & \texttt{7340} & R & \texttt{} & \texttt{ST   R1, [R5]} & {\footnotesize DISPLAY = 3 : running again} \\
\texttt{1A} & \texttt{F000} & N & \texttt{} & \texttt{HALT} & {\footnotesize } \\

\end{listtab}
\end{document}
